% ********************************************************************
% Parameters file - Edit this to customize your thesis information
% ********************************************************************

% ====================== THESIS METADATA =========================

% Your name
\newcommand{\myName}{Leon Håkansson Ingman}

% Thesis title
\newcommand{\myTitle}{Machine Learning Detection of Shadow Fleet Vessels in the Baltic Sea}

% Degree
\newcommand{\myDegree}{Master of Science in Engineering}

% Program
\newcommand{\myProgram}{Computer Science and Engineering}

% Department
\newcommand{\myDepartment}{Department of Computer Science}

% Faculty
\newcommand{\myFaculty}{Faculty of Engineering, LTH}

% University
\newcommand{\myUni}{Lund University}

% Submission date
\newcommand{\myTime}{June 2025}
\newcommand{\myYear}{2025}

% Supervisors (replace with actual names when known)
\newcommand{\mySupervisor}{[Supervisor Name], [Title]}
\newcommand{\myCoSupervisor}{} % Leave empty if no co-supervisor

% Examiner (replace with actual name when known)
\newcommand{\myExaminer}{[Examiner Name], [Title]}

% Company/Organization (if external thesis)
\newcommand{\myCompany}{Rendavouz}

% Abstract (English) - Edit this with your actual abstract
\newcommand{\myAbstract}{%
This thesis investigates the application of machine learning techniques for detecting shadow fleet vessels in the Baltic Sea. Shadow fleet vessels, which operate outside normal maritime regulations and tracking systems, pose significant challenges for maritime security and environmental protection. 

Using vessel tracking data and machine learning algorithms, this work develops a predictive model capable of identifying vessels with characteristics typical of shadow fleet operations. The methodology involves pattern recognition of vessel routes, anomaly detection in behavior, and classification based on known shadow fleet vessel features.

[Add more details about your specific approach, datasets, results, and contributions]

\textbf{Keywords:} Machine Learning, Shadow Fleet, Vessel Detection, Maritime Security, Baltic Sea, AIS Data, Anomaly Detection
}

% Acknowledgements (optional - edit as needed)
\newcommand{\myAcknowledgements}{%
I would like to thank [Supervisor Name] for guidance and support throughout this thesis work. I am grateful to my partners at Rendavouz for providing the opportunity to work on this important problem and for access to data and resources.

Special thanks to [add other people you want to thank].

This work was carried out during the spring semester of 2025 at the Department of Computer Science, Lund University.
}

% Popular summary in Swedish - Edit this with your actual popular summary
\newcommand{\mySammanfattning}{%
[Skriv din populärvetenskapliga sammanfattning på svenska här. Detta ska vara en sammanfattning av ditt examensarbete som är begriplig för en bredare publik utan teknisk bakgrund. Förklara problemet, din metod och dina resultat på ett tillgängligt sätt.]

Exempel:
Skuggflottan är fartyg som opererar utanför normala maritima regleringar och spårningssystem. Dessa fartyg utgör betydande utmaningar för sjöfartssäkerhet och miljöskydd i Östersjön. Detta examensarbete undersöker hur maskininlärning kan användas för att identifiera skuggflottans fartyg baserat på deras rörelsemönster och beteende.

[Fortsätt med din faktiska sammanfattning...]
}
